\documentclass[11pt]{article}

\usepackage[margin=1in]{geometry}
\usepackage{amsmath,amssymb,booktabs,longtable,array}
\usepackage{hyperref}
\usepackage{microtype}

\title{Robert's Feedback as Formula Proposals}
\author{}
\date{}

\newcommand{\sg}{\operatorname{sg}}
\newcommand{\ProjectMeanOne}{\operatorname{ProjectMeanOne}}
\newcommand{\E}{\mathbb{E}}
\newcommand{\Var}{\operatorname{Var}}
\newcommand{\Cov}{\operatorname{Cov}}
\newcommand{\tr}{\operatorname{tr}}

\begin{document}
\maketitle

This note translates Robert's comments into mathematical variants that can be
folded into the AD-CISPO / CW-CISPO proposal.

\section*{Glossary}

\begin{longtable}{>{\raggedright\arraybackslash}p{0.18\linewidth}p{0.30\linewidth}p{0.44\linewidth}}
\toprule
Symbol & Meaning & Notes \\
\midrule
\endhead
$q$ & Prompt or question & Fixed input before generation starts. \\
$i$ & Response index & One sampled response within a group of completions for the same prompt. \\
$G$ & Group size & Number of sampled responses for one prompt. \\
$o_i$ & Sampled response $i$ & Generated token sequence. \\
$o_{i,t}$ & Token at position $t$ in response $i$ & A generated token, not a prompt token, when $t \in M_i$. \\
$T_i$ & Length of response $i$ & Number of generated-token positions before masking. \\
$t$ & Target token position & Position whose saliency, credit weight, or clip multiplier is being computed. \\
$j$ & Future token position & Used when a later token attends back to position $t$. \\
$u$ & Dummy summation index & Used when averaging over all valid positions in $M_i$. \\
$s_{i,t}$ & Autoregressive state before token $o_{i,t}$ & Prompt plus generated prefix before position $t$. \\
$M_i$ & Valid generated-token action positions for response $i$ & Excludes prompt tokens, padding, masked-out actions, special tokens, and sink tokens removed by the saliency guard. All token-level means, sums, and mean-one normalizations are over this set. \\
$\pi_\theta$ & Current policy being optimized & The model receiving the policy-gradient update. \\
$\pi_{\theta_{\mathrm{old}}}$ & Rollout policy & The policy that produced the sampled batch. \\
$\rho_{i,t}$ & Token importance-sampling ratio & Current-policy probability divided by rollout-policy probability for token $o_{i,t}$. \\
$\hat A_i$ & Sequence-level advantage & Usually group-relative for the sampled response. \\
$r_i$ & Outcome reward for response $i$ & Usually a non-differentiable verifier reward. \\
$F$ & Teacher-forced scalar objective & Current causal-tangent proxy differentiates this objective, not reward directly. \\
$h^{(\ell)}_{i,t}$ & Hidden state & Layer $\ell$, response $i$, position $t$. \\
$S_{i,t}$ & Generic token saliency score & Sections below define attention, causal-tangent, rollout, entropy, and opportunity variants. \\
$H_{i,t}$ & Token entropy & Uncertainty in the next-token distribution. \\
$C_{i,t}$ & Controllability score & How much action choice at a state changes downstream return. \\
$w_{i,t}$ & CW-CISPO token-credit weight & Detached, bounded, and normalized to mean one over $M_i$. \\
$m_{i,t}$ & AD-CISPO clip multiplier & Used to redistribute the upper clipping radius. \\
$P_i$ & Token attention transition matrix & Used for PageRank-style multi-step attention saliency. \\
$c_i$ & PageRank-style token centrality vector & Fixed-point centrality over generated-token positions. \\
\bottomrule
\end{longtable}

\section*{Shared Notation}

For prompt $q$, let response index $i$ select one sampled completion:
\begin{equation}
  o_i = (o_{i,1}, \ldots, o_{i,T_i}).
\end{equation}

Let
\begin{equation}
  s_{i,t} = (q, o_{i,<t})
\end{equation}
be the autoregressive state before token $o_{i,t}$.

The set $M_i$ is the set of valid generated-token action positions for response
$i$. It is the denominator set for token-level means and the support for credit
weights. In practice, $M_i$ excludes prompt positions, padding positions,
masked-out generated actions, special tokens, and separator or punctuation sink
tokens removed by the saliency guard. Thus $M_i$ is the set of generated
positions eligible to receive saliency mass, CW-CISPO credit weight, or
AD-CISPO clipping budget.

The token importance-sampling ratio is
\begin{equation}
  \rho_{i,t}(\theta)
  =
  \frac{
    \pi_\theta(o_{i,t} \mid s_{i,t})
  }{
    \pi_{\theta_{\mathrm{old}}}(o_{i,t} \mid s_{i,t})
  }.
\end{equation}

\section*{1. Current Proxy: Teacher-Forced Objective Sensitivity}

The current causal-tangent saliency is approximately a hidden-state sensitivity
to a teacher-forced scalar objective $F$, not directly to the outcome reward.

One clean way to write the current objective is
\begin{equation}
  F_i(\theta)
  =
  |\hat A_i|
  \frac{1}{|M_i|}
  \sum_{t \in M_i}
  \log \pi_\theta(o_{i,t} \mid s_{i,t}),
\end{equation}
with batch objective
\begin{equation}
  F(\theta)
  =
  \sum_i F_i(\theta).
\end{equation}

The current causal-tangent score is then
\begin{equation}
  S^{F}_{i,t}
  =
  \frac{1}{|L_{\mathrm{top}}|}
  \sum_{\ell \in L_{\mathrm{top}}}
  \frac{
    \left|
      \nabla_{h^{(\ell)}_{i,t}} F
      \cdot
      \left(
        h^{(\ell)}_{i,t}
        -
        \frac{1}{|M_i|}
        \sum_{u \in M_i} h^{(\ell)}_{i,u}
      \right)
    \right|
  }{\sqrt{d_{\mathrm{model}}}}.
\end{equation}

In words:
\begin{equation}
  S^{F}_{i,t}
  \approx
  \text{local effect of changing } h_{i,t}
  \text{ on teacher-forced objective } F.
\end{equation}

Robert's point is that the desired object is closer to reward sensitivity:
\begin{equation}
  S^{r}_{i,t}
  \approx
  \left\|
    \nabla_{h_{i,t}} r_i
  \right\|.
\end{equation}

Since $r_i$ is usually a non-differentiable verifier reward, the differentiable
ideal is more like
\begin{equation}
  S^{r}_{i,t}
  \approx
  \left\|
    \nabla_{h_{i,t}}
    \E_{o_{i,>t} \sim \pi_\theta}
    \left[
      r(q, o_{i,\le t}, o_{i,>t})
    \right]
  \right\|.
\end{equation}

\section*{2. Rollout-Based Reward Impact, VinePPO-Style}

The expensive but conceptually clean target is the counterfactual effect of
token $o_{i,t}$ on final reward.

Define the continuation value after prefix $o_{i,\le t}$:
\begin{equation}
  Q^\pi(s_{i,t}, o_{i,t})
  =
  \E_{o_{i,>t} \sim \pi}
  \left[
    r(q, o_{i,\le t}, o_{i,>t})
  \right].
\end{equation}

Define the pre-token value
\begin{equation}
  V^\pi(s_{i,t})
  =
  \E_{a \sim \pi(\cdot \mid s_{i,t})}
  \left[
    Q^\pi(s_{i,t}, a)
  \right].
\end{equation}

Then a reward-impact saliency is
\begin{equation}
  S^{\mathrm{rollout}}_{i,t}
  =
  \left|
    Q^\pi(s_{i,t}, o_{i,t})
    -
    V^\pi(s_{i,t})
  \right|.
\end{equation}

With $K$ sampled continuations per candidate token, estimate
\begin{equation}
  \widehat Q^\pi(s_{i,t}, a)
  =
  \frac{1}{K}
  \sum_{k=1}^{K}
  r(q, o_{i,<t}, a, o^{(k)}_{>t}).
\end{equation}

The rollout saliency estimator becomes
\begin{equation}
  \widehat S^{\mathrm{rollout}}_{i,t}
  =
  \left|
    \widehat Q^\pi(s_{i,t}, o_{i,t})
    -
    \sum_{a}
    \pi(a \mid s_{i,t})
    \widehat Q^\pi(s_{i,t}, a)
  \right|.
\end{equation}

\section*{3. Learned Offline Predictor for Grad-H $F$ or Reward Impact}

Instead of computing expensive saliency online every batch, train a small
predictor $g_\phi$ from hidden states and cheap token features.

For vector-gradient prediction:
\begin{equation}
  g_\phi(h^{(\ell)}_{i,t}, x_{i,t})
  \approx
  \sg
  \left(
    \nabla_{h^{(\ell)}_{i,t}} F
  \right),
\end{equation}
with loss
\begin{equation}
  \mathcal{L}_{\mathrm{grad}}(\phi)
  =
  \E_{i,t,\ell}
  \left[
    \left\|
      g_\phi(h^{(\ell)}_{i,t}, x_{i,t})
      -
      \sg
      \left(
        \nabla_{h^{(\ell)}_{i,t}} F
      \right)
    \right\|_2^2
  \right].
\end{equation}

For scalar saliency prediction:
\begin{equation}
  \hat S^\phi_{i,t}
  =
  g_\phi(h_{i,t}, x_{i,t})
  \approx
  S^F_{i,t}
  \quad \text{or} \quad
  S^r_{i,t}.
\end{equation}

Train with
\begin{equation}
  \mathcal{L}_{\mathrm{score}}(\phi)
  =
  \E_{i,t}
  \left[
    \left(
      g_\phi(h_{i,t}, x_{i,t})
      -
      \sg(S^\star_{i,t})
    \right)^2
  \right],
\end{equation}
where
\begin{equation}
  S^\star_{i,t}
  \in
  \left\{
    S^F_{i,t},
    S^r_{i,t},
    S^{\mathrm{rollout}}_{i,t}
  \right\}.
\end{equation}

Then CW-CISPO can use
\begin{equation}
  w_{i,t}
  =
  \ProjectMeanOne
  \left(
    \left(
      \hat S^\phi_{i,t} + \epsilon
    \right)^\gamma
  \right).
\end{equation}

\section*{4. Probability-Space Sensitivity}

Robert was unsure about saliency written as $\nabla_{p(o_i)} F$. The concern is
that perturbing the probability of a sampled token is not the same as
perturbing a hidden state or a policy parameter.

Let
\begin{equation}
  p_{i,t}
  =
  \pi_\theta(o_{i,t} \mid s_{i,t}).
\end{equation}

A probability-coordinate sensitivity would be
\begin{equation}
  S^p_{i,t}
  =
  \left|
    p_{i,t}
    \frac{\partial F}{\partial p_{i,t}}
  \right|.
\end{equation}

Equivalently, because
\begin{equation}
  \frac{\partial F}{\partial \log p_{i,t}}
  =
  p_{i,t}
  \frac{\partial F}{\partial p_{i,t}},
\end{equation}
we can write
\begin{equation}
  S^p_{i,t}
  =
  \left|
    \frac{\partial F}{\partial \log \pi_\theta(o_{i,t} \mid s_{i,t})}
  \right|.
\end{equation}

For the simple teacher-forced objective above,
\begin{equation}
  \frac{\partial F}{\partial \log \pi_\theta(o_{i,t} \mid s_{i,t})}
  =
  \frac{|\hat A_i|}{|M_i|}.
\end{equation}
That is nearly uniform within a response, so it is not a useful token saliency
unless $F$ includes richer token-level structure.

\section*{5. Decision-State and Bottleneck-State Formulations}

In deep RL terms, a token position is important if the next action materially
changes the reachable return distribution.

A decision-state score is
\begin{equation}
  D_{i,t}
  =
  \Var_{a \sim \pi(\cdot \mid s_{i,t})}
  \left[
    Q^\pi(s_{i,t}, a)
  \right].
\end{equation}

Equivalently, a max-gap version is
\begin{equation}
  D^{\max}_{i,t}
  =
  \max_a Q^\pi(s_{i,t}, a)
  -
  \min_a Q^\pi(s_{i,t}, a).
\end{equation}

A bottleneck-state score can be written as the association between visiting a
state-like feature $z$ and eventual success:
\begin{equation}
  B(z)
  =
  \Pr(z \in \tau \mid r=1)
  -
  \Pr(z \in \tau \mid r=0).
\end{equation}

For token positions, define a feature map $\psi(s_{i,t})$. Then
\begin{equation}
  B_{i,t}
  =
  B(\psi(s_{i,t})).
\end{equation}

The doorway analogy from grid-world RL corresponds to states with both high
visitation centrality and high success association:
\begin{equation}
  S^{\mathrm{bottleneck}}_{i,t}
  =
  C^{\mathrm{visit}}_{i,t}
  \cdot
  \left|B_{i,t}\right|.
\end{equation}

\section*{6. Multi-Step Attention Centrality, PageRank-Style}

Current future-attention in-degree is a one-step score:
\begin{equation}
  S^{\mathrm{attn},1}_{i,t}
  =
  \sum_{\ell,h}
  \sum_{\substack{j \in M_i \\ j > t}}
  a^{(\ell,h)}_{i,j,t}.
\end{equation}
This implicitly weights every future query token $j$ equally.

Define a row-stochastic token transition matrix $P_i$ over generated positions:
\begin{equation}
  P_i[j,t]
  =
  \frac{1}{|L_{\mathrm{top}}|}
  \sum_{\ell \in L_{\mathrm{top}}}
  \frac{1}{H_\ell}
  \sum_{h=1}^{H_\ell}
  a^{(\ell,h)}_{i,j,t}.
\end{equation}

A PageRank-style equilibrium centrality $c_i$ satisfies
\begin{equation}
  c_i
  =
  (1-\alpha) \nu_i
  +
  \alpha P_i^\top c_i,
  \qquad
  0 < \alpha < 1,
\end{equation}
where $\nu_i$ can be uniform over $M_i$:
\begin{equation}
  \nu_i[t]
  =
  \frac{1}{|M_i|}.
\end{equation}

The multi-step attention saliency is
\begin{equation}
  S^{\mathrm{attn},\mathrm{PR}}_{i,t}
  =
  c_i[t].
\end{equation}

Equivalently,
\begin{equation}
  c_i
  =
  (1-\alpha)
  \sum_{n=0}^{\infty}
  \alpha^n
  \left(P_i^\top\right)^n
  \nu_i.
\end{equation}

This is attention paid by future tokens, weighted by those tokens' own
centrality.

\section*{7. Multiplicative Saliency, Entropy, and Controllability}

The current saliency-only form is
\begin{equation}
  \tilde w_{i,t}
  =
  \left(S_{i,t} + \epsilon\right)^\gamma.
\end{equation}

Robert's suggestion is closer to a multiplicative gate: update tokens only
when they are both important and changeable.

Let token entropy be
\begin{equation}
  H_{i,t}
  =
  -
  \sum_a
  \pi_\theta(a \mid s_{i,t})
  \log \pi_\theta(a \mid s_{i,t}).
\end{equation}

Let controllability be
\begin{equation}
  C_{i,t}
  =
  \Var_{a \sim \pi_\theta(\cdot \mid s_{i,t})}
  \left[
    Q^\pi(s_{i,t}, a)
  \right].
\end{equation}

A simple multiplicative weight is
\begin{equation}
  \tilde w^{\mathrm{mult}}_{i,t}
  =
  \left(S_{i,t} + \epsilon\right)^\gamma
  \left(H_{i,t} + \epsilon\right)^\eta
  \left(C_{i,t} + \epsilon\right)^\kappa.
\end{equation}

Then enforce bounded mean-one credit weights:
\begin{equation}
  w^{\mathrm{mult}}_{i,t}
  =
  \ProjectMeanOne
  \left(
    \tilde w^{\mathrm{mult}}_{i,t};
    w_{\min},
    w_{\max}
  \right),
\end{equation}
so that
\begin{equation}
  \frac{1}{|M_i|}
  \sum_{t \in M_i}
  w^{\mathrm{mult}}_{i,t}
  =
  1.
\end{equation}

An Osband-style delight gate can be written as a sigmoid:
\begin{equation}
  g_{i,t}
  =
  \sigma
  \left(
    \beta_0
    +
    \beta_{\ell}
    \left[-\log \pi_\theta(o_{i,t} \mid s_{i,t})\right]
    +
    \beta_A \hat A_i
    +
    \beta_C \log(C_{i,t}+\epsilon)
  \right).
\end{equation}

Then
\begin{equation}
  \tilde w^{\mathrm{gate}}_{i,t}
  =
  \left(S_{i,t}+\epsilon\right)^\gamma
  g_{i,t}.
\end{equation}

The intuition is
\begin{equation}
  \text{credit}
  \approx
  \text{importance}
  \times
  \text{uncertainty}
  \times
  \text{controllability}.
\end{equation}

\section*{8. Cumulative Future Entropy Baseline}

A cheap baseline is cumulative future entropy: tokens are important if many
uncertain choices remain downstream.

Define entropy-to-go:
\begin{equation}
  E^{\mathrm{future}}_{i,t}
  =
  \sum_{\substack{j \in M_i \\ j \ge t}}
  \lambda^{j-t}
  H_{i,j},
  \qquad
  0 < \lambda \le 1.
\end{equation}

The corresponding mean-one credit weight is
\begin{equation}
  w^{\mathrm{futureH}}_{i,t}
  =
  \ProjectMeanOne
  \left(
    \left(E^{\mathrm{future}}_{i,t}+\epsilon\right)^\eta
  \right).
\end{equation}

\section*{9. Token-Dependent Learning Rate View}

CW-CISPO weights can also be read as a token-dependent learning rate.

The scalar CISPO gradient has the form
\begin{equation}
  g_{\mathrm{CISPO}}
  =
  -
  \sum_{i,t}
  \sg(\hat \rho_{i,t})
  \hat A_i
  \nabla_\theta
  \log \pi_\theta(o_{i,t} \mid s_{i,t}).
\end{equation}

CW-CISPO changes it to
\begin{equation}
  g_{\mathrm{CW}}
  =
  -
  \sum_{i,t}
  \sg(w_{i,t})
  \sg(\hat \rho_{i,t})
  \hat A_i
  \nabla_\theta
  \log \pi_\theta(o_{i,t} \mid s_{i,t}).
\end{equation}

This is equivalent to a per-token step size
\begin{equation}
  \eta_{i,t}
  =
  \eta_0 w_{i,t}.
\end{equation}

Mean-one normalization preserves the average learning rate:
\begin{equation}
  \frac{1}{|M_i|}
  \sum_{t \in M_i}
  \eta_{i,t}
  =
  \eta_0.
\end{equation}

Robert's learning-rate point can be stated as an implicit relative-step decay:
\begin{equation}
  \frac{\|\Delta \theta\|_2}{\|\theta\|_2}
  \approx
  \eta
  \frac{\|g\|_2}{\|\theta\|_2}.
\end{equation}

Even with constant $\eta$, this relative step shrinks as $\|\theta\|_2$ grows.
Token-dependent weights can selectively counteract or amplify this decay at
important states.

\section*{10. Feature-Space Opportunity Proxy}

Robert's "opportunities opened by this token" idea can be formalized as
successor-feature dispersion.

Let $\phi(q,o_{\le t},o_{>t})$ be a feature vector of the completed trajectory,
such as final hidden state, answer embedding, or verifier-side features. Define
the feature opportunity of prefix $o_{\le t}$:
\begin{equation}
  O_{i,t}
  =
  \tr
  \left(
    \Cov_{o_{>t} \sim \pi}
    \left[
      \phi(q,o_{\le t},o_{>t})
    \right]
  \right).
\end{equation}

An action-conditional opportunity gain is
\begin{equation}
  \Delta O_{i,t}
  =
  O(q,o_{\le t})
  -
  \E_{a \sim \pi(\cdot \mid s_{i,t})}
  \left[
    O(q,o_{<t},a)
  \right].
\end{equation}

Use
\begin{equation}
  S^{\mathrm{opp}}_{i,t}
  =
  \left|\Delta O_{i,t}\right|
\end{equation}
as a proxy for whether token $o_{i,t}$ opens or closes useful future trajectory
options.

\section*{11. Golden Trajectory, Teacher Logits, and PRM Access}

There are three increasingly rich supervision signals.

\paragraph{Golden trajectory only.}
Given only a teacher or gold trajectory $o^T$, the objective is
\begin{equation}
  F_{\mathrm{gold}}(\theta)
  =
  \sum_{t}
  \log \pi_\theta(o^T_t \mid q,o^T_{<t}).
\end{equation}
The saliency is
\begin{equation}
  S^{\mathrm{gold}}_t
  =
  \left\|
    \nabla_{h_t} F_{\mathrm{gold}}
  \right\|.
\end{equation}

\paragraph{Teacher logits.}
If teacher logits or a teacher distribution $\tau_t(a)$ are available, use
\begin{equation}
  F_{\mathrm{logits}}(\theta)
  =
  \sum_t
  \sum_a
  \tau_t(a)
  \log \pi_\theta(a \mid q,o^T_{<t}).
\end{equation}

Equivalently,
\begin{equation}
  F_{\mathrm{logits}}(\theta)
  =
  -
  \sum_t
  \operatorname{KL}
  \left(
    \tau_t(\cdot)
    \,\|\, 
    \pi_\theta(\cdot \mid q,o^T_{<t})
  \right)
  + \mathrm{const}.
\end{equation}

The saliency is
\begin{equation}
  S^{\mathrm{logits}}_t
  =
  \left\|
    \nabla_{h_t} F_{\mathrm{logits}}
  \right\|.
\end{equation}

\paragraph{Process reward model.}
If a process reward model $R_\psi$ is available, define
\begin{equation}
  F_{\mathrm{PRM}}(\theta)
  =
  \sum_t
  R_\psi(q,o_{\le t}).
\end{equation}
The saliency is
\begin{equation}
  S^{\mathrm{PRM}}_t
  =
  \left\|
    \nabla_{h_t} F_{\mathrm{PRM}}
  \right\|.
\end{equation}

The hierarchy is therefore
\begin{equation}
  \text{gold trajectory}
  \;<\;
  \text{teacher logits}
  \;\not\equiv\;
  \text{process reward model}.
\end{equation}

\section*{12. Cost of Full-Trajectory Grad-H $F$}

The current causal-tangent formulation requires a full teacher-forced forward
pass over the sampled trajectory:
\begin{equation}
  (q,o_i)
  \longmapsto
  \{h^{(\ell)}_{i,t}\}_{\ell \in L_{\mathrm{top}},\, t \in M_i}
  \longmapsto
  F_i
  \longmapsto
  \nabla_{h^{(\ell)}_{i,t}} F_i.
\end{equation}

The cost scales as
\begin{equation}
  \operatorname{cost}
  \left(
    S^F
  \right)
  \approx
  \operatorname{Forward}(T_i)
  +
  \operatorname{Backward}_{h}(T_i, |L_{\mathrm{top}}|),
\end{equation}
with activation storage
\begin{equation}
  \operatorname{memory}
  \left(
    S^F
  \right)
  =
  O
  \left(
    |L_{\mathrm{top}}| \, T_i \, d_{\mathrm{model}}
  \right).
\end{equation}

This is why cheap baselines are important:
\begin{equation}
  S^{\mathrm{attn},1}_{i,t},
  \qquad
  S^{\mathrm{attn},\mathrm{PR}}_{i,t},
  \qquad
  S^{\mathrm{futureH}}_{i,t},
  \qquad
  \hat S^\phi_{i,t}.
\end{equation}

\section*{13. Baseline Matrix}

The concrete baseline family can be summarized as
\begin{align}
  S^{\mathrm{attn},1}_{i,t}
  &=
  \sum_{\ell,h}
  \sum_{j>t}
  a^{(\ell,h)}_{i,j,t},
  \\
  S^{\mathrm{attn},\mathrm{PR}}_{i,t}
  &=
  c_i[t],
  \quad
  c_i=(1-\alpha)\nu_i+\alpha P_i^\top c_i,
  \\
  S^F_{i,t}
  &=
  \left\|
    \nabla_{h_{i,t}} F
  \right\|_{\mathrm{tangent}},
  \\
  S^r_{i,t}
  &=
  \left\|
    \nabla_{h_{i,t}}
    \E[r \mid q,o_{\le t}]
  \right\|,
  \\
  S^{\mathrm{rollout}}_{i,t}
  &=
  \left|
    Q^\pi(s_{i,t},o_{i,t})
    -
    V^\pi(s_{i,t})
  \right|,
  \\
  S^{\mathrm{futureH}}_{i,t}
  &=
  \sum_{j \ge t}
  \lambda^{j-t}
  H_{i,j},
  \\
  S^{\mathrm{opp}}_{i,t}
  &=
  \left|
    \Delta O_{i,t}
  \right|.
\end{align}

Each score becomes a CISPO credit weight through the same bounded mean-one map:
\begin{equation}
  w_{i,t}
  =
  \ProjectMeanOne
  \left(
    (S_{i,t}+\epsilon)^\gamma;
    w_{\min},
    w_{\max}
  \right).
\end{equation}

The main experimental question is:
\begin{equation}
  \text{Does } w(S^F),\; w(S^r),\; \text{or } w(S^{\mathrm{rollout}})
  \text{ beat }
  w(S^{\mathrm{futureH}})
  \text{ and shuffled controls?}
\end{equation}

\end{document}
